\section{Hankel matrices and finitely many inverse limits}
\label{sec:hankel}

We construct a positive semidefinite Hankel matrix of rank $20$ and a
$20$-dimensional subspace $W$ with a finiteness property: certain nonzero
limits of inverses of compressions to $W$ have only finitely many
projective directions. The construction has two parts. A rank calculation,
proved in Appendix~\ref{app:certificate}, supplies a point at which a
kernel-projection map has sufficiently large derivative. We then divide
out translations and dilations, choose a finite fiber of the resulting
map, and use that fiber to classify the inverse limits.

\subsection{The rank stratum and the seed matrix}

Let $x=(x_0,x_1,\ldots,x_7)$ and put
\[
 H=\R[x]_2=L\oplus P,
 \qquad
 L=\langle x_0^2,x_0x_1,\ldots,x_0x_7\rangle,
 \qquad
 P=\R[x_1,\ldots,x_7]_2.
\]
Subscripts denote homogeneous degree. Monomials are our coordinate bases,
and these bases are orthonormal whenever we take an orthogonal complement
or compression. Setting $x_0=1$ identifies $H$ with polynomials of degree
at most two in $s=(s_1,\ldots,s_7)$. In this identification, $L$ consists
of the constant and linear terms and $P$ consists of the quadratic terms.
Write
\begin{equation}
\label{eq:hankel-dimensions}
 h=\dim H=36,\qquad d_*=20,\qquad k=h-d_*=16,
 \qquad q=d_*-\dim L=12.
\end{equation}

Let $\mathcal H\subset\Sym(H)$ be the space of Hankel matrices
\begin{equation}
\label{eq:hankel-functional}
 Y(p,p')=y(pp'),\qquad p,p'\in H,
 \qquad y\in\R[x]_4^*.
\end{equation}
Multiplication $\Sym^2H\to\R[x]_4$ is onto, so $y\mapsto Y$ is
injective and $\dim\mathcal H=\binom{11}{4}=330$.
We also call $Y$ a moment matrix. This terminology refers only to
\eqref{eq:hankel-functional}; no representing measure is assumed.
In the chart $x_0=1$, the functional $y$ acts on polynomials in $s$ of
degree at most four.

For subspaces of polynomial spaces, a product such as $JH_0$ denotes
the linear span of all products of one element of each space. We use
the following algebraic input.

\begin{lemma}[Seed matrix]
\label{lem:seed}
There is a positive semidefinite matrix $T\in\mathcal H$ of rank $d_*=20$
with the following properties. Its restriction to $L$ is positive
definite. If
\[
 J=\ker T,\qquad K=\operatorname{proj}_P J,
 \qquad H_0=\{p\in H:T(p,L)=0\},
\]
then $\dim K=k=16$, multiplication
\[
 \Sym^2K\longrightarrow\R[x_1,\ldots,x_7]_4
\]
is injective, and
\begin{equation}
\label{eq:pure-cubic-spanning}
 K\langle x_1,\ldots,x_7\rangle
   =\R[x_1,\ldots,x_7]_3.
\end{equation}
Finally, $JH_0$ has codimension $9$ in $\R[x]_4$.
\end{lemma}

The proof in Appendix~\ref{app:certificate} begins with two cyclic
orbits of real points and deforms their moment matrix inside the
rank-$20$ stratum. All rank conditions needed for that deformation are
verified there by explicit minors.

Let $\mathcal X$ be the set of projective classes $[Y]\in\mathbf P(\mathcal H)$
such that $\rank Y=d_*$, $\ker Y\cap L=0$, and the projected kernel
$K_Y=\operatorname{proj}_P(\ker Y)$ satisfies the two multiplication
conditions of Lemma~\ref{lem:seed}. Positivity is not part of the
definition of $\mathcal X$. Define
\[
 \pi:\mathcal X\longrightarrow\Gr(k,P),
 \qquad \pi([Y])=K_Y.
\]

\begin{lemma}
\label{lem:hankel-stratum}
The set $\mathcal X$ is a smooth semialgebraic manifold of dimension
\[
 D=330-1-\frac{k(k+1)}2=193.
\]
Before projectivization, the tangent space at a representative $Y$ is
the space of $Y'\in\mathcal H$ whose restriction to
$\ker Y\times\ker Y$ is zero.
\end{lemma}

\begin{proof}
Put $J_Y=\ker Y$. Taking highest homogeneous parts in the affine chart
identifies $J_Y$ with $K_Y$. Injectivity of multiplication on
$\Sym^2K_Y$ therefore implies injectivity on $\Sym^2J_Y$: a relation
among products in $J_Y$ would have a relation among their highest parts.
By duality, restriction of Hankel forms gives a surjection
\[
 \mathcal H\longrightarrow\Sym(J_Y).
\]
Choose a complement of $J_Y$ in $H$. The compression of $Y$ to this
complement is invertible: a vector in the complement orthogonal for
$Y$ to that complement is orthogonal for $Y$ to all of $H$, and hence
lies in $J_Y$. In a neighborhood of $Y$, rank $d_*$ is thus equivalent
to the vanishing of a $k$-by-$k$ symmetric Schur complement. At $Y$,
the derivative of this Schur complement is restriction to $J_Y$.
The implicit function theorem gives a smooth manifold of codimension
$k(k+1)/2=136$ in $\mathcal H$, with the asserted tangent space.
The condition $\ker Y\cap L=0$ and the multiplication conditions on
$K_Y$ are open rank conditions. Finally, projectivization removes the
nonzero scalar parameter, giving dimension $193$.
\end{proof}

\subsection{Removing translations and dilations}

The fibers of $\pi$ contain the orbits of affine translations and
positive dilations in the $s$ variables. We give explicit coordinates
for removing these eight parameters. Because $J_Y\cap L=0$, its
elements, in the affine chart, have a unique description
\begin{equation}
\label{eq:kernel-graph}
 J_Y=\{p(s)+b(p)(s)+a(p):p\in K_Y\},
\end{equation}
where $b:K_Y\to\R[s]_1$ takes values in homogeneous linear forms and
$a:K_Y\to\R$ is linear. These coefficient maps depend continuously
and semialgebraically on $[Y]$.

Pushforward of moments by $s\mapsto s+t$ transforms a kernel
polynomial by $f(s)\mapsto f(s-t)$. Thus translation changes the
coefficient maps by
\begin{equation}
\label{eq:translation-kernel-coefficients}
 b_t(p)=b(p)-\partial_t p,
 \qquad
 a_t(p)=a(p)-b(p)(t)+p(t),
\end{equation}
where $\partial_t p$ is the directional derivative of $p$ in direction
$t$. Pushforward by $s\mapsto\lambda s$, with $\lambda>0$, gives
\begin{equation}
\label{eq:dilation-kernel-coefficients}
 b\longmapsto\lambda b,
 \qquad a\longmapsto\lambda^2 a.
\end{equation}
Both actions are congruences on moment matrices. They preserve Hankel
structure, rank, and $\pi$, and preserve positive semidefiniteness when
it is present.

\begin{lemma}
\label{lem:kernel-affine-action}
For every $[Y]\in\mathcal X$, the map
\[
 \Delta_{K_Y}:\R^7\longrightarrow
   \operatorname{Hom}(K_Y,\R[s]_1),
 \qquad t\longmapsto(p\mapsto\partial_t p)
\]
is injective, and the coefficient maps $a,b$ in
\eqref{eq:kernel-graph} are not both zero. The seven infinitesimal
translations and the infinitesimal dilation give eight independent
variations of the kernel graph. These assertions, including the
independence of the variations, also hold over $\mathbf C$ under the
same rank and multiplication hypotheses.
\end{lemma}

\begin{proof}
If $t\ne0$ belongs to the kernel of $\Delta_{K_Y}$, a linear change of
the pure variables makes $t$ the first coordinate direction. All
quadratics in $K_Y$ are then independent of that coordinate, so their
products with linear forms cannot span its cube. This contradicts
\eqref{eq:pure-cubic-spanning}.

Suppose next that $a=b=0$, so $J_Y=K_Y$ in the chart. The kernel
condition and \eqref{eq:pure-cubic-spanning} force all moments of
degree three to vanish. Multiplying that spanning identity by linear
forms shows that $K_YP=\R[s]_4$, so all moments of degree four vanish
as well. With coordinates ordered by degrees $0,1,2$, the moment
matrix consequently has the form
\[
 \begin{pmatrix}
  *&*&*\\ *&*&0\\ *&0&0
 \end{pmatrix},
\]
where the degree-zero block has size one and the degree-one block
has size seven. The last block row has rank at most one, and the
first eight rows have rank at most eight. Thus $\rank Y\le9$, a
contradiction.

For independence of the infinitesimal variations, suppose a translation
in direction $t$ and a dilation with infinitesimal coefficient $\alpha$
have zero combined variation of the graph. Equations
\eqref{eq:translation-kernel-coefficients} and
\eqref{eq:dilation-kernel-coefficients} give
\[
 \alpha b(p)=\partial_t p,
 \qquad 2\alpha a(p)=b(p)(t)
 \quad(p\in K_Y).
\]
If $\alpha=0$, injectivity of $\Delta_{K_Y}$ gives $t=0$.
If $\alpha\ne0$, put $c=t/\alpha$. Then
$b(p)=\partial_c p$ and $a(p)=p(c)$, so every polynomial in the graph
is $p(s+c)$. A translation would therefore produce the already
excluded case $a=b=0$. This proves independence. Each argument is
algebraic and applies equally over $\mathbf C$.
\end{proof}

Equip the spaces of coefficient maps with the Euclidean inner
products induced by the monomial bases. We normalize $[Y]$ in two
steps. First choose the unique translation for which
$b\perp\ran\Delta_{K_Y}$. Explicitly, its parameter is
\[
 t=(\Delta_{K_Y}^*\Delta_{K_Y})^{-1}\Delta_{K_Y}^*b.
\]
Next choose the unique positive dilation for which
\begin{equation}
\label{eq:kernel-normalization}
 \|b\|^2+\|a\|^2=1.
\end{equation}
For the translated coefficients, this amounts to solving
$\lambda^2\|b\|^2+\lambda^4\|a\|^2=1$. The left side is strictly
increasing from zero to infinity, by Lemma~\ref{lem:kernel-affine-action}.

\begin{lemma}
\label{lem:hankel-normalization}
The normalized matrices, together with their data $(K_Y,b,a)$, form
a semialgebraic set $\mathcal M$ of dimension $185$. Normalization is
continuous, and each $[Y]\in\mathcal X$ is uniquely recovered from
its normalized data and its seven translation parameters and positive
dilation parameter. Moreover, $\mathcal M$ has a compact ambient
semialgebraic space in which the projection to $\Gr(k,P)$ is
continuous.
\end{lemma}

\begin{proof}
The translation formula is continuous because $\Delta_{K_Y}$ has
constant full column rank. The unique positive solution for $\lambda$
is continuous and semialgebraic as well. Dilation preserves the
orthogonality condition on $b$, so the two normalizations are compatible.
Conversely, start with normalized data, undo its positive dilation,
and then undo its translation. Applying the two normalization steps
recovers the prescribed parameters. This gives a semialgebraic
homeomorphism between $\mathcal X$ and
$\mathcal M\times\R^7\times\R_{>0}$, and hence
$\dim\mathcal M=193-8=185$.

To specify a compact ambient space, retain the projective matrix
$[Y]$ and the subspace $K_Y$, and extend $a,b$ by zero on $K_Y^\perp$
in $P$. The extended maps have the same norms as the original maps.
Condition \eqref{eq:kernel-normalization} places their pair on a
fixed unit sphere. The product of this sphere with
$\mathbf P(\mathcal H)$ and $\Gr(k,P)$ is compact, and projection
to the last factor is continuous. We take the closure of
$\mathcal M$ in this product.
\end{proof}

\subsection{Selecting a finite fiber}

We now show that the positive semidefinite part of $\mathcal X$ has
enough distinct projected kernels to avoid both infinite normalized
fibers and limits at the boundary of $\mathcal M$. We use the standard
dimension and fiber-dimension theorems for semialgebraic sets, as well
as the fact that the frontier $\overline S\setminus S$ of a
semialgebraic set has smaller dimension than $S$; see
\cite{BCR1998}. A zero-dimensional semialgebraic set is finite.

At any $T$ representing a point of $\mathcal X$, put $J=\ker T$ and
$H_0=\{p:T(p,L)=0\}$. The kernel of the derivative of the
kernel-projection map, before projectivizing, is
\begin{equation}
\label{eq:projection-tangent}
 \{Y'\in\mathcal H:Y'(J)\subset T(L)\}
       =(JH_0)^\perp.
\end{equation}
Here the right side is an annihilator in $\R[x]_4^*$, identified
with $\mathcal H$. To verify the identity, continue a vector
$j\in J$ along a differentiable curve of kernels. Its first variation
$j'$ must satisfy
\[
 Y'j+Tj'=0.
\]
The projected kernel has zero first variation precisely when the
pure part of each $j'$ lies in $K$. Subtracting an element of $J$
then makes $j'$ belong to $L$, without changing this equation.
Thus the condition is $Y'J\subset T(L)$. Conversely, that condition
allows all the corrections to be chosen in $L$; it also implies
$Y'|_{J\times J}=0$, so Lemma~\ref{lem:hankel-stratum} gives an
actual rank-stratum tangent. Finally,
$H_0=T(L)^\perp$, which proves the annihilator description.

At the seed matrix of Lemma~\ref{lem:seed}, the space in
\eqref{eq:projection-tangent} has dimension $9$. The rank stratum
before projectivizing has dimension $194$, so the derivative of
$\pi$ has rank $194-9=185$. Projectivization removes one kernel
direction, namely scaling, and leaves this derivative rank unchanged.
Near the seed matrix, rank remains $20$ and its nonzero eigenvalues
remain positive. Its $L$ block remains positive definite as well.
Consequently the $\pi$-image of matrices with these positivity
properties has dimension at least $185$. It has dimension at most
$185$ because $\pi$ factors through $\mathcal M$.

The seed's final condition can also be retained in this neighborhood.
Indeed, an invertible $321$-row product minor for $JH_0$ remains
invertible nearby. On the other hand, scaling, translations, and
dilation always give nine independent tangent directions in
\eqref{eq:projection-tangent}: the eight affine directions have
independent kernel variations by Lemma~\ref{lem:kernel-affine-action},
whereas scaling has zero kernel variation. Thus $\dim(JH_0)\le321$
throughout this neighborhood.

\begin{proposition}
\label{prop:finite-normalized-fiber}
The matrix $T$ in Lemma~\ref{lem:seed} may be chosen so that, for
$K=\operatorname{proj}_P\ker T$,
\begin{enumerate}
 \item the fiber $\mathcal M_K$ of $\mathcal M\to\Gr(k,P)$ is finite;
 \item no point of $\overline{\mathcal M}\setminus\mathcal M$ projects
 to $K$.
\end{enumerate}
\end{proposition}

\begin{proof}
The set of subspaces over which the fiber of $\mathcal M$ has positive
dimension has dimension at most $184$, by the fiber-dimension theorem
and $\dim\mathcal M=185$. The frontier of $\mathcal M$ has dimension
at most $184$, so its projection to $\Gr(k,P)$ has dimension at most
$184$ as well. The positive semidefinite image just constructed has
dimension $185$. We can therefore choose $K$ in that image outside
both exceptional sets, and choose a corresponding positive
semidefinite $T$ in the neighborhood under consideration. Its fiber
is nonempty and zero-dimensional, hence finite.
\end{proof}

Fix this $T$ and $K$ for the rest of the paper, and set
\begin{equation}
\label{eq:fixed-complements}
 R=K^\perp\subset P,\qquad W=L\oplus R.
\end{equation}
Thus $H=W\oplus K$ is an orthogonal decomposition, $\dim R=q=12$,
and $\dim W=d_*=20$.

Each normalized class $[\bar Y]\in\mathcal M_K$ determines a
nonzero symmetric matrix on $W$, up to scale. Let $\bar A$ be the
compression of a representative $\bar Y$ to $W$. The kernel of $\bar Y$
is a graph over $K$ with lower terms in $L$, so $W$ is a complement of
that kernel. The same nondegeneracy argument used in
Lemma~\ref{lem:hankel-stratum} shows that $\bar A$ is invertible,
even if $\bar Y$ is indefinite. Define $C_{\bar Y}\in\Sym(W)$ to
have block form
\begin{equation}
\label{eq:finite-inverse-directions}
 C_{\bar Y}=
 \begin{pmatrix}0&0\\0&(\bar A^{-1})_{R,R}\end{pmatrix}
 \quad\hbox{on }W=L\oplus R,
 \qquad
 \mathcal C=\{[C_{\bar Y}]:[\bar Y]\in\mathcal M_K\}.
\end{equation}
If its $R,R$ block were zero, the invertible operator $\bar A^{-1}$
would map $R$ into $L$, which is impossible because $12>8$.
Thus $\mathcal C$ is a finite nonempty set of projective classes of
nonzero symmetric matrices on $W$.

\subsection{The inverse-limit property}

To connect the normalized fiber to arbitrary limiting sequences, we
first need a rank correction that remains small even when some
matrix entries diverge. All norms below are fixed Euclidean operator
norms; convergence is equivalent in any choice of such norms.

\begin{lemma}[Uniform rank correction]
\label{lem:hankel-rank-correction}
Use the fixed decomposition $H=W\oplus K$ from
\eqref{eq:fixed-complements}. Suppose $Y_i\in\mathcal H$ have blocks
\[
 Y_i=\begin{pmatrix}A_i&B_i\\B_i^T&D_i\end{pmatrix},
\]
where $A_i$ is invertible, $\sup_i\|A_i^{-1}\|<\infty$, and
\[
 A_i^{-1}B_i\longrightarrow0,
 \qquad D_i-B_i^TA_i^{-1}B_i\longrightarrow0.
\]
There are Hankel perturbations $\Delta_i\to0$ such that
$Y_i^\sharp=Y_i+\Delta_i$ has rank $d_*$. Its $W$ compression
$A_i^\sharp$ satisfies
\[
 (A_i^\sharp)^{-1}-A_i^{-1}\longrightarrow0,
\]
and $\ker Y_i^\sharp\to K$, or equivalently
$\ran Y_i^\sharp\to W$.
\end{lemma}

\begin{proof}
On matrices whose $W$ block is invertible, write
$\mathcal S(Y)=D-B^TA^{-1}B$ for the Schur complement on $K$.
For a variation $\Delta$ with blocks $(a,b,d)$ and $F=A^{-1}B$,
\begin{equation}
\label{eq:schur-differential}
 D\mathcal S_Y(\Delta)
   =d-b^TF-F^Tb+F^TaF.
\end{equation}
At $Y_i$, these derivatives tend to the restriction map
$\Delta\mapsto\Delta_{K,K}$ on Hankel space. That restriction map
is onto $\Sym(K)$ because multiplication on $\Sym^2K$ is injective.
Choose a fixed linear right inverse
$\mathcal E:\Sym(K)\to\mathcal H$.

We record why an implicit-function argument is uniform in $i$.
Put $M=\sup_i\|A_i^{-1}\|$. On a fixed sufficiently small ball of
perturbations, the inverse of the perturbed $W$ block obeys
\[
 (A_i+a)^{-1}=(I+A_i^{-1}a)^{-1}A_i^{-1},
 \qquad \|(A_i+a)^{-1}\|\le2M.
\]
Moreover,
\[
 (A_i+a)^{-1}(B_i+b)
   =(I+A_i^{-1}a)^{-1}(A_i^{-1}B_i+A_i^{-1}b)
\]
is uniformly bounded there. Differentiating the inverse and this
last expression gives
\[
 \delta(A^{-1})=-A^{-1}(\delta A)A^{-1},
 \qquad
 \delta F=A^{-1}(\delta B-(\delta A)F).
\]
Together with \eqref{eq:schur-differential}, these identities give a
uniform bound for the second derivatives of $\mathcal S$ on that
ball. This bound uses neither boundedness of $A_i$ nor boundedness
of $B_i$.

Define $f_i(z)=\mathcal S(Y_i+\mathcal E z)$ for
$z\in\Sym(K)$ in a fixed small ball. We have $f_i(0)\to0$ and
$Df_i(0)\to I$. Shrinking the fixed ball if needed, the uniform
second-derivative bound gives
$\|Df_i(z)-I\|\le1/2$ there for all sufficiently large $i$.
The map $z\mapsto z-f_i(z)$ is therefore a contraction on the
closed ball of radius $2\|f_i(0)\|$: its value at zero has norm
$\|f_i(0)\|$, and its Lipschitz constant is at most $1/2$.
Its fixed point $z_i$ satisfies $f_i(z_i)=0$ and $z_i\to0$.
Set $\Delta_i=\mathcal E z_i$.

The $W$ block remains invertible, and its Schur complement is zero,
so $\rank Y_i^\sharp=d_*$. The inverse identity gives
$(A_i^\sharp)^{-1}-A_i^{-1}\to0$. The displayed formula for the
perturbed cross block also gives
$(A_i^\sharp)^{-1}B_i^\sharp\to0$. The kernel of $Y_i^\sharp$ is
the graph with columns
$\binom{-(A_i^\sharp)^{-1}B_i^\sharp}{I_K}$, proving the final
assertion. Discarding finitely many indices has no effect on any
of the conclusions.
\end{proof}

\begin{lemma}[Finite inverse limits]
\label{lem:inverse-limits}
Let $T,K,R,W$, and $\mathcal C$ be chosen as above. Suppose a
sequence $Y_i\in\mathcal H$ has a group of $d_*$ eigenvalues bounded
away from zero in absolute value, and the span of the corresponding
eigenvectors tends to $W$. Suppose all remaining eigenvalues tend
to zero. If the inverses of the compressions $A_i=(Y_i)_{W,W}$
converge to a nonzero matrix $C$, then
\[
 [C]\in\mathcal C.
\]
\end{lemma}

\begin{proof}
We first put the sequence in the situation of
Lemma~\ref{lem:hankel-rank-correction}. Split off the distinguished
eigenvalues to write $Y_i=Y_i^{\mathrm{large}}+E_i$, where
$\rank Y_i^{\mathrm{large}}=d_*$ and $E_i\to0$. After passage to a
subsequence, orthonormal frames for the large eigenspaces have
$W$ components $S_i$ converging to an orthogonal matrix and $K$
components tending to zero. In particular $S_i$ is invertible, and
\begin{equation}
\label{eq:spectral-graph}
 Y_i^{\mathrm{large}}=
 \begin{pmatrix}I\\G_i\end{pmatrix}
 A_i^0
 \begin{pmatrix}I&G_i^T\end{pmatrix},
 \qquad G_i\longrightarrow0,
 \qquad \sup_i\|(A_i^0)^{-1}\|<\infty.
\end{equation}
Indeed $A_i^0=S_i\operatorname{diag}(\lambda_{i,1},\ldots,
\lambda_{i,d_*})S_i^T$, and the inverse eigenvalues are bounded.

Write the blocks of $E_i$ as $(a_i,b_i,d_i)$ in the same order as
in the rank-correction lemma. Then
\[
 A_i=A_i^0+a_i,
 \qquad B_i=A_i^0G_i^T+b_i,
 \qquad D_i=G_iA_i^0G_i^T+d_i.
\]
The inverses $A_i^{-1}$ are bounded, and
\[
 A_i^{-1}B_i=G_i^T+A_i^{-1}(b_i-a_iG_i^T)\longrightarrow0.
\]
With $e_i=b_i-a_iG_i^T$, direct block elimination gives
\[
 D_i-B_i^TA_i^{-1}B_i
 =d_i-G_ib_i-b_i^TG_i^T+G_ia_iG_i^T-e_i^TA_i^{-1}e_i
 \longrightarrow0.
\]
These formulas explain explicitly why divergence of the large
eigenvalues causes no difficulty. Apply
Lemma~\ref{lem:hankel-rank-correction}. The corrected Hankel matrices
have exact rank $d_*$, their kernels tend to $K$, and their
compression inverses have the same nonzero limit $C$. We replace
$Y_i$ by these corrected matrices for the rest of the proof.

For large $i$, the class $[Y_i]$ belongs to $\mathcal X$: its kernel
is close to $K\subset H$, and the required product conditions are
open at $K$. Put $K_i=\pi([Y_i])$. Then $K_i\to K$, and the lower
coefficient maps $a_i^{\mathrm{ker}},b_i^{\mathrm{ker}}$ of the kernel
graph tend to zero.
The translations used in normalization consequently tend to zero,
by the explicit left-inverse formula for $\Delta_{K_i}$. The
translated coefficient maps still tend to zero, so the normalizing
dilation factors tend to infinity. Let their reciprocals be
$\eps_i\to0$.

The normalized data lie in the compact ambient space of
Lemma~\ref{lem:hankel-normalization}. Pass to a convergent
subsequence. Its limit projects to $K$, and
Proposition~\ref{prop:finite-normalized-fiber} excludes a limit in
the frontier of $\mathcal M$. It is therefore a member
$[\bar Y]\in\mathcal M_K$. Choose representatives $Z_i$ of the
normalized matrices converging to a representative $\bar Y$.
Up to nonzero scalar multiples, the original matrices are obtained
from $Z_i$ by dilation by $\eps_i$, followed by translations whose
parameters tend to zero.

We compute the inverse limit first before those translations.
Set $R_i=K_i^\perp\subset P$ and $W_i=L\oplus R_i$. Use orthonormal
coordinates on $W_i$ converging to the fixed coordinates on $W$,
and keep the coordinates on $L$ unchanged. The compressions
$\bar A_i=(Z_i)_{W_i,W_i}$ then converge to the invertible
compression $\bar A$ of $\bar Y$ to $W$. Dilation acts on monomial
columns by
\[
 D_{\eps}=\diag(1,\eps I_7,\eps^2 I_{28})
\]
on $H$, according to degrees $0,1,2$. Since it preserves $W_i$,
the compression of the dilated matrix is
\[
 A_i^{\mathrm{dil}}
   =D_{\eps_i,W_i}\bar A_iD_{\eps_i,W_i}.
\]
Consequently
\begin{equation}
\label{eq:dilated-inverse-limit}
 \eps_i^4(A_i^{\mathrm{dil}})^{-1}
 =\diag(\eps_i^2,\eps_i I_7,I_q)\,
   \bar A_i^{-1}\,
   \diag(\eps_i^2,\eps_i I_7,I_q)
 \longrightarrow C_{\bar Y}.
\end{equation}

It remains to justify changing $W_i$ to $W$ and restoring the
translations. After dilation, the kernel coefficients of $Z_i$
are $\eps_i b_i^{\mathrm{norm}}$ and
$\eps_i^2 a_i^{\mathrm{norm}}$. The normalized coefficients are
bounded. Write the dilated kernel in $W_i\oplus K_i$ as
$\{(F_i k,k):k\in K_i\}$; the preceding coefficient bounds give
$F_i\to0$. Its orthogonal complement, which is the range of the
dilated matrix, has graph embedding
\[
 \iota_i:W_i\longrightarrow H,\qquad
 \iota_i w=(w,-F_i^*w).
\]
Thus $\iota_i$ tends to the inclusion of $W$ after the chosen
coordinate identifications, and the full dilated matrix is
$\iota_iA_i^{\mathrm{dil}}\iota_i^T$.

Let $\mathsf B_i$ denote the monomial matrix of the remaining translation.
Its parameter tends to zero, so $\mathsf B_i\to I_H$. Compression to $W$
after this translation therefore gives
\[
 M_iA_i^{\mathrm{dil}}M_i^T,
 \qquad M_i=\operatorname{proj}_W\mathsf B_i\iota_i\longrightarrow I_W.
\]
The actual compression $A_i$ differs from this matrix by a nonzero
scalar $\alpha_i$, coming from the choice of projective
representative. Equation~\eqref{eq:dilated-inverse-limit} yields
\[
 \alpha_i\eps_i^4 A_i^{-1}
   =M_i^{-T}\bigl(\eps_i^4(A_i^{\mathrm{dil}})^{-1}\bigr)M_i^{-1}
   \longrightarrow C_{\bar Y}\ne0.
\]
Hence $[A_i^{-1}]\to[C_{\bar Y}]$. Since $A_i^{-1}\to C\ne0$
also gives $[A_i^{-1}]\to[C]$, we obtain
$[C]=[C_{\bar Y}]\in\mathcal C$.
\end{proof}

We retain the positive semidefinite matrix $T$, the subspaces
$K$, $R=K^\perp\subset P$, and $W=L\oplus R$, and the finite set
$\mathcal C$ from \eqref{eq:finite-inverse-directions}. These are the
geometric data used in the next construction.
